\subsection{OpenCL -- Open Computing Language}

OpenCL:
\begin{itemize}
    \item \enquote{OpenCL is a trademark of Apple Inc. used under license to the Khronos Group Inc}
    \item working group created 2008~\cite{khronos_opencl_working_group_khronos_2008}
    \item first speciifcation published 2009~\cite{munshi_opencl_2009}
    \item current standard OpenCL 3.0 revision 17 October 24, 2024~\cite{khronos_opencl_working_group_opencl_2024}
    \item Implementations OpenCL 3.0 (\url{https://www.khronos.org/conformance/adopters/conformant-products/opencl})
    \begin{itemize}
        \item Intel
        \item Imagination Technologies
        \item QUALCOMM
        \item Google, Inc. 
        \item Software Freedom Conservancy
        \item Arm Limited
        \item NVIDIA
        \item Software in the Public Interest, Inc.
        \item Verisilicon, Inc.
        \item pocl
    \end{itemize}
    \item many more for OpenCL 2.0 (like AMD, Samsung Electronics, etc.)
\end{itemize}

\begin{itemize}
    \item \acrfull{opencl}
    \item standardized
    \item kernels are simple strings $\rightarrow$ easy to manipulate, tedious to use
    \item kernels must be compiled by yourself using OpenCL provided functions
    \item C API
    \item \gls{jit} compiled
    \item programming model/kernel structure close to \gls{cuda}
\end{itemize}

\begin{advantagebox}
    \begin{itemize}
        \item standardized
        \item supports a wide variety of hardware platforms (many implementations)
    \end{itemize}
\end{advantagebox}

\begin{disadvantagebox}
    \begin{itemize}
        \item kernels must be compiled by hand
        \item error prone
        \item some low-level optimizations not possible
    \end{itemize}
\end{disadvantagebox}


\begin{listing}
    \begin{minted}{cpp}
const char* kernel_source = R"KS(
__kernel void saxpy(const double alpha, __global double *X, __global double *X) {
    const int idx = get_global_id(0);
    Y[idx] = alpha * X[idx] + Y[idx];
}
)KS";

cl_platform_id platform;
clGetPlatformIDs(1, &platform, nullptr);
cl_device_id device;
clGetDeviceIDs(platform, CL_DEVICE_TYPE_GPU, 1, &device, nullptr);
cl_context context = clCreateContext(nullptr, 1, &device, nullptr, nullptr, nullptr);
cl_command_queue queue = clCreateCommandQueue(context, device, 0, nullptr);

cl_mem d_X = clCreateBuffer(context, CL_MEM_READ_ONLY | CL_MEM_COPY_HOST_PTR, N * sizeof(double), X, nullptr);
cl_mem d_Y = clCreateBuffer(context, CL_MEM_READ_ONLY | CL_MEM_COPY_HOST_PTR, N * sizeof(double), Y, nullptr);

cl_program program = clCreateProgramWithSource(context, 1, &kernel_source, nullptr, nullptr);
clBuildProgram(program, 1, &device, nullptr, nullptr, nullptr);
cl_kernel kernel = clCreateKernel(program, "saxpy", nullptr);

clSetKernelArg(kernel, 0, sizeof(double), &alpha);
clSetKernelArg(kernel, 1, sizeof(cl_mem), &d_X);
clSetKernelArg(kernel, 2, sizeof(cl_mem), &d_Y);

size_t global_size = N;
clEnqueueNDRangeKernel(queue, kernel, 1, nullptr, &global_size, nullptr, 0, nullptr, nullptr);
clEnqueueReadBuffer(queue, d_C, CL_TRUE, 0, N * sizeof(double), C, 0, nullptr, nullptr);

clReleaseMemObject(d_X); 
clReleaseMemObject(d_Y);
clReleaseKernel(kernel); 
clReleaseProgram(program);
clReleaseCommandQueue(queue); 
clReleaseContext(context);
    \end{minted}
    \caption{%
    SAXPY example using \gls{opencl}.
    }
    \label{code:opencl_example}
\end{listing}